% p3-11-limitations

\section{P3 §11. Limitations}

11. Limitations

   Figure (fig-p3-ch11-limitations). Limitations triptych (Paper III) --- Independence open / 42 not a
                                         proof set / Coq status.

  11.1 Unproven Independence of Generators
  Propositions 3.1 and 3.5, and Corollary 3.6, are all conditional on the Q-linear
  independence of log 3, log $\varphi$, log pi, 1. Baker's theorem on linear forms in logarithms of
  algebraic numbers covers log 3 and log $\varphi$ (both algebraic), but logpi involves
  transcendental arguments not covered by the standard machinery. The status of logpi
  relative to Q-linear independence of algebraic logarithms is an open problem in
  transcendence theory. All results depending on this independence are marked
  "(conditional)" throughout.

  Consequence. If independence fails --- i.e., if there exists a non-trivial Q-linear relation
  among log 3, log $\varphi$, log pi, 1 --- then the Trinity monomial lattice T is not a free abelian
  group in log-space, and collisions can occur: two distinct labels may represent the same
  real number. In this case, the MDL score would need to be adjusted to account for the
  effective dimensionality of the lattice.

  11.2 AlphaPhi Notation (Resolved for Paper 3)
  The notation inconsistency between alpha$\varphi$( alg) approx 0.118 and alpha$\varphi$(log) approx
  0.306 in the early drafts of Chapter 38 of the PhD program is resolved by Section 10.5
  above. Paper 3 uses the disambiguated notation throughout. All future papers in this
  trilogy must adopt this disambiguation; Section 10.5 is the canonical reference.

  11.3 Floating-Point Limitations
  All numerical examples use 50-digit approximations. Computed Pellis coefficients and
  Trinity approximants are valid only to the working precision. Physical constants are
  given to current best-known values, with CODATA 2022 used for NIST constants; as
  measurements improve, the expansions change. The framework is applied to current

  best values, not exact transcendental numbers, and all numerical results carry implicit
  uncertainty at the 50th decimal digit.

  11.4 The 42-Entry Catalogue Is Not a Proof Set
  The 42 precision fits of the PhD program's Chapter 34 are an empirical survey, not a
  proof that the listed symbolic forms are the structurally "correct" representations of the
  constants. The MDL reinterpretation (Section 7.3) preserves this status and explicitly
  labels each entry as a compression hypothesis, not a theoretical derivation. High-MDL
  catalogue entries provide evidence against the represented form, not evidence for
  anything else.

  The corresponding Coq file t27/proofs/trinity/Catalog42.v should be read as a
  representative formal-provenance layer for the catalogue. In the available source audit,
  that file contains 13 source-level Qed. markers, 0 Admitted. markers, and a
  totalverifiedtheorems aggregate over 28 representative theorem obligations. The
  surrounding t27/proofs/trinity/ layer is larger: 13 Coq files, 281 theorem-like
  declarations, 122 source-level Qed. markers, and 32 source-level Admitted. markers.
  This is stronger than an informal table, but it is not yet a one-row-one-theorem
  formalization of all 42 catalogue entries.

  11.5 Physical Constants Are Not Explained
  Compact Trinity representations of alpha QED-1 or mupe are purely numerical facts
  about current measured values. The MDL score measures description economy, not
  causal or theoretical relationships. No claim is made about the physical origin of any
  constant.

  11.6 Non-Uniqueness of Optimal Approximants
  At any complexity budget L, two Trinity monomials may be equidistant from the target.
  The argmin in Definition 5.1 is generically unique (under the working independence
  hypothesis) but can fail for special targets. A canonical tiebreaking rule is adopted in
  Appendix A.3 (lexicographic on label); this rule is implemented in the reference
  software.

  11.7 Computational Feasibility
  The Trinity lattice search is naively O(2L) in the complexity budget L. Practical
  feasibility requires L $\leq$ 25 without logarithmic pruning; with the log-space pruning of
  Algorithm 9.2, this extends to L $\leq$ 40 or so. Theorem 7.2 (MDL optimality) applies only
  to the exact minimiser; the pruned algorithm is a heuristic approximation.

  11.8 Conjecture Status of Core Results

        Conjecture 6. 5 (spectral gap) and Conjecture 8.1 (phase transition) are open.
        Their proof would require development of Baladi-type functional-analytic
        machinery in the GOLDEN BRIDGE (Trinity-Pellis Bridge) setting and, for the phase transition, a connection
        to homogeneous dynamics (Ratner theory) that remains to be formalised.

  11.9 Coq Proof Status
  Identities from the PhD program's Appendix F (Coq citation map) are cited as formal
  provenance, not as complete proofs accessible to general mathematical verification.
  Exact theorem counts are intentionally not reproduced here because the Coq audit
  ledger is regenerated as the repository evolves. Paper 3 uses only local status labels:
  Qed, Admitted / proof-sketch, not formalised, or open. Any theorem whose dependency
  chain touches an Admitted item must be treated as conditional until the regenerated
  ledger closes it.

\subsection{References}

- Vasilev, D. \textit{Flos Aureus Monograph: Trinity Constants}. DOI \href{https://zenodo.org/doi/10.5281/zenodo.19227877}{10.5281/zenodo.19227877}.
- CODATA 2018 recommended values: NIST \href{https://physics.nist.gov/cuu/Constants/}{https://physics.nist.gov/cuu/Constants/}.
- Heyrovska, R. Golden-angle FSC publications --- Heyrovský Institute of Physical Chemistry.
